\section{The list that became a connection}
\label{sec:ch05-the-list-that-became-a-connection}
\index{pagination!nested connection|(}

Chapter~\ref{ch:hotchocolate-quickly} shipped \code{Product.reviews} as
\code{[Review!]!} and said in as many words that a later chapter would make it
a connection. This is that chapter, and the field is the second thing here a
client cannot survive.

\begin{graphqlsdl}
reviews(
  first: Int
  after: String
  last: Int
  before: String
): ReviewConnection!
\end{graphqlsdl}

There is no additive path to this. \code{reviews} was a list and is now an
object with \code{nodes}, \code{edges}, \code{pageInfo} and
\code{totalCount}~\autocite{relay2026connections}; every client that selected it
breaks. The alternative was to leave \code{reviews} alone and add a second field
beside it, which is what chapter~\ref{ch:data-without-the-n-plus-1} did with
\code{browseProducts}, and I did not do it twice for a reason worth stating.
\code{products} had two chapters of published measurements riding on it, so
keeping it was buying something specific. \code{reviews} has nothing riding on
it except this book's own tests, and a schema that grows a second field every
time the first one is wrong ends up as two schemas. Take the breakage while
there is nobody to break.

\subsection*{The attribute that does not do what it says}

\index{HotChocolate!UseConnection}%
My first attempt put the paging attribute on the resolver, returned
\code{Connection<Review>}, compiled clean, and produced this:

\begin{graphqlsdl}
reviews: [Review!]! @cost(weight: "10")
\end{graphqlsdl}

A list. No error, no warning, no diagnostic. The field I had annotated as a
connection was still a list, and the only reason I noticed is that I export the
schema and read it, which is the habit
chapter~\ref{ch:hotchocolate-quickly} built for exactly this class of surprise.
No tag holds that state, because it was a wrong turn rather than a step, but the
lab rebuilds it in three lines if you want to watch it happen.

\code{[UseConnection]} does not make a field a connection. Its own summary in
the source says what it is for, and it is not what the name suggests: it
overrides the global paging options for a specific field. What it adds is a
validation middleware and a set of options. It never touches the field's type.
The attribute that rewrites the type is \code{[UsePaging]}, whose configuration
calls \code{UsePaging} on the field descriptor.

What actually made \code{reviews} a connection is the return type:

\begin{minted}[fontsize=\footnotesize]{csharp}
[UseConnection(IncludeTotalCount = true)]
public static async Task<PageConnection<Review>> GetReviewsAsync(
    [Parent("Id")] Product product,
    IReviewsByProductIdDataLoader reviewsByProductId,
    PagingArguments pagingArguments,
    CancellationToken cancellationToken)
{
    var page = await reviewsByProductId
        .With(pagingArguments)
        .LoadAsync(product.Id, cancellationToken);

    return page ?? Page<Review>.Empty;
}
\end{minted}

\code{PageConnection<TNode>} carries a name attribute that expands to the node
type followed by \code{Connection}, and declares an implicit conversion from
\code{Page<TNode>}, which is why the last line of that method returns the wrong
type and compiles anyway. The attribute is tuning; the return type is the
contract.

I am spending a page on a mistake I made because the failure mode is the
chapter's second silent one, and it is worse than
chapter~\ref{ch:data-without-the-n-plus-1}'s. That one produced an all-zero
identifier at runtime, which at least happens. This one produces a schema that
is quietly not the schema you thought you wrote, and the only detector is
reading the SDL.

\subsection*{One statement for twenty-five pages}

\index{pagination!ToBatchPageAsync}%
A connection on a child field is not the same problem as a connection on a root
field, and this is where the chapter earns its keep. \code{browseProducts} pages
one list. \code{reviews} pages twenty-five lists, one per product on the page,
and every one of them needs its own first three rows.

Fetching every review for every product and slicing in memory would make the
connection a lie, because the whole argument for paging a child collection is
that the rows nobody asked for never leave the database.
\code{ToBatchPageAsync} is what makes the honest version one statement:

\begin{minted}[fontsize=\footnotesize]{csharp}
var pages = await db.Reviews
    .AsNoTracking()
    .Where(r => productIds.Contains(r.ProductId))
    .OrderBy(r => r.CreatedAt)
    .ThenBy(r => r.Id)
    .ToBatchPageAsync(r => r.ProductId, pagingArguments, cancellationToken);
\end{minted}

It returns a dictionary of pages, one per parent key. The \code{OrderBy} is not
decoration: the cursor keys are read off the ordering already on the queryable,
and the tiebreaker on the identifier is what makes that ordering total, for the
same reason chapter~\ref{ch:data-without-the-n-plus-1} gave for
\code{browseProducts}.

Here is what PostgreSQL is asked, for \code{reviews(first: 3)}:

\begin{minted}[fontsize=\footnotesize]{sql}
SELECT r1.product_id, r3.id, r3.body, r3.created_at, r3.customer_id,
       r3.product_id, r3.rating
FROM (
    SELECT r.product_id FROM reviews AS r
    WHERE r.product_id = ANY (@productIds)
    GROUP BY r.product_id
) AS r1
LEFT JOIN (
    SELECT r2.id, r2.body, r2.created_at, r2.customer_id, r2.product_id, r2.rating
    FROM (
        SELECT r0.id, r0.body, r0.created_at, r0.customer_id, r0.product_id,
               r0.rating,
               ROW_NUMBER() OVER(PARTITION BY r0.product_id
                                 ORDER BY r0.created_at, r0.id) AS row
        FROM reviews AS r0
        WHERE r0.product_id = ANY (@productIds)
    ) AS r2
    WHERE r2.row <= 4
) AS r3 ON r1.product_id = r3.product_id
ORDER BY r1.product_id, r3.product_id, r3.created_at, r3.id
\end{minted}

\index{PostgreSQL!window function}%
A window function, partitioned by the parent key. PostgreSQL numbers each
product's reviews in the sort order and throws away everything past the cut, so
twenty-five separate pages come back in one round trip.

Read the cut carefully: \code{row <= 4}, for a page of three. The extra row is
how \code{hasNextPage} is answered. A connection has to say whether there is
more, and the cheapest honest way to know is to ask for one more than you intend
to return and see whether it arrives. If you have ever wondered why a page of
ten sometimes reads eleven rows, that is why.

\subsection*{Why the DataLoader has to split in two}

\index{DataLoader!branching}%
The DataLoader takes \code{PagingArguments} as a parameter, and it is not an
ordinary service parameter. The generator recognises the type and emits a
\code{GetRequiredState} call for it, which means the arguments arrive as
DataLoader \emph{state} rather than as part of the key. What puts them there is
the \code{.With} call in the resolver, and that call does something else at the
same time:

\begin{minted}[fontsize=\footnotesize]{csharp}
var branchKey = pagingArguments.ComputeHash(context);
var state = new PagingState<TValue>(pagingArguments, context);
return (IQueryDataLoader<TKey, Page<TValue>>)dataLoader.Branch(
    branchKey, CreateBranch, state);
\end{minted}

It hashes the page shape and hands back a \emph{branch} of the DataLoader keyed
on that hash. A branch wraps the same underlying loader with its own state and,
crucially, its own promise cache.

\index{DataLoader!cache key}%
The necessity is easier to see from the failure it prevents. A DataLoader caches
by key, and the key here is a product identifier. Two fields in one request
asking for two different pages of the same product's reviews would both look up
the same key, and one of them would be handed the other's page: correct-looking
data, wrong rows, no error. Branching makes the page shape part of the identity
of the loader rather than of the key, so the two cannot see each other at all.
Figure~\ref{fig:ch05-branching} draws the case I measured: one product, asked
for two different pages of its own reviews in a single document.

\begin{figure}[htbp]
  \centering
  % Why a paged DataLoader has to branch. Referenced from chapter 05. Bare
% tikzpicture: the figure environment, caption and label live at the call site.
% Uses only arrows.meta and positioning. The numbers are the measured ones from
% the aliased two-page query the lab asks for.
\begin{tikzpicture}[
    font=\small,
    field/.style={draw, rounded corners=2pt, minimum width=34mm,
                  minimum height=8mm, align=center, font=\scriptsize},
    branch/.style={draw, fill=black!8, rounded corners=2pt, minimum width=34mm,
                   minimum height=9mm, align=center, font=\scriptsize},
    base/.style={draw, thick, rounded corners=2pt, minimum width=44mm,
                 minimum height=7mm, align=center, font=\scriptsize},
    key/.style={-{Stealth[length=1.6mm]}, gray},
    fat/.style={-{Stealth[length=2mm]}, thick},
    note/.style={font=\scriptsize\itshape, align=center}
  ]

  % ---- one generated DataLoader ----------------------------------------
  \node[base] (gen) at (0,1.5) {ReviewsByProductIdDataLoader};

  % ---- two call sites asking for different pages -----------------------
  \node[field] (f1) at (-3.3,-0.2) {\texttt{a: reviews(first: 2)}};
  \node[field] (f2) at (3.3,-0.2)  {\texttt{b: reviews(first: 5)}};

  \draw[key] (f1.north) -- (gen.south west);
  \draw[key] (f2.north) -- (gen.south east);

  \node[note, text width=38mm] at (0,0.0)
    {\texttt{.With(...)} hashes\\the page shape};

  % ---- two branches ----------------------------------------------------
  \node[branch] (b1) at (-3.3,-2.0) {branch \texttt{f:2}\\own cache, own batch};
  \node[branch] (b2) at (3.3,-2.0)  {branch \texttt{f:5}\\own cache, own batch};

  \draw[fat] (f1) -- (b1);
  \draw[fat] (f2) -- (b2);

  \draw[key] (-3.3,-2.5) -- (-3.3,-3.1);
  \draw[key] (3.3,-2.5) -- (3.3,-3.1);

  \node[field] (s1) at (-3.3,-3.6) {1 statement\\\texttt{row <= 3}};
  \node[field] (s2) at (3.3,-3.6)  {1 statement\\\texttt{row <= 6}};

  \node[note, text width=68mm] at (0,-5.0)
    {One branch would have to serve both, and a promise cache keyed
     only on the product would hand the second alias the first
     alias's page.};

\end{tikzpicture}

  \caption{One generated DataLoader, two page shapes, two branches. Everything
  asking for the same page shares a branch and therefore a batch, which is what
  keeps twenty-five products at one statement. A different page shape gets its
  own branch and its own statement, and the two window cuts in the database log
  show which is which.}
  \label{fig:ch05-branching}
\end{figure}

That document costs three statements: one for the product, and one for each
branch. The cuts in the two of them are \code{row <= 3} and \code{row <= 6},
which is the n+1 rule from the last subsection applied to a page of two and a
page of five separately. Nothing coordinated them, and nothing had to.

The hash covers \code{first}, \code{after}, \code{last} and \code{before}, and
also the filter, sort and projection when a \code{QueryContext} is passed. That
is a wider notion of identity than the batching in
chapter~\ref{ch:data-without-the-n-plus-1} needed, and it is the right one:
anything that changes what a fetch would return has to change which loader you
are talking to.

\subsection*{An empty answer, again}

\index{DataLoader!missing key}%
\code{ToBatchPageAsync} returns entries only for keys that had rows, and three
of Mosaic's twenty-five products have never been reviewed. \code{reviews} is
non-nullable. That is exactly the shape of
chapter~\ref{ch:data-without-the-n-plus-1}'s trap, arriving again in a new
costume.

The escape used last time is not available. That chapter fixed it by returning
an \code{ILookup}, which answers an unknown key with an empty sequence, and a
page is not a sequence. So the empty answer has to be written down:

\begin{minted}[fontsize=\footnotesize]{csharp}
foreach (var productId in productIds)
{
    pages.TryAdd(productId, Page<Review>.Empty);
}
\end{minted}

Two lines, and without them three products answer null for a non-nullable field
and the request fails. What I take from meeting this twice is that the choice of
return type on a batch fetch is a decision about semantics rather than about
convenience, and the question to ask each time is what this thing says when it
has nothing to say.

\subsection*{The numbers, and the one that went the wrong way}

\index{pagination!measured cost}%
The chapter's query has to be written differently now:

\begin{minted}[fontsize=\footnotesize]{graphql}
{ products { title
    reviews(first: 12) { nodes { rating author { displayName } } } } }
\end{minted}

\code{first: 12} is not arbitrary. The most-reviewed product has exactly twelve,
so this still asks for all hundred and twenty reviews, and the number the
verification script asserts is still the number chapters two to four asserted.

\begin{minted}[fontsize=\footnotesize]{text}
parse - validate - compile - coerce - execute 5.405ms total 5.529ms
    (document cache hit, operation cache hit, 146 resolvers, 3 SQL)
\end{minted}

A hundred and forty-six resolvers and three statements. The same two numbers
chapter~\ref{ch:data-without-the-n-plus-1} finished on, through a change that
replaced the field's type, its resolver, its DataLoader and the SQL underneath
it. Ask for a smaller page and the resolver count falls to 88, because 88 is one
for \code{products} plus twenty-five for \code{reviews} plus sixty-two authors,
and sixty-two is what you get by asking twenty-five products for three reviews
each when some of them have fewer. The statement count stays at three. The
window returns fewer rows; it does not run more times.

Now the number that did not flatter the change, and one caveat about the
milliseconds in the line above: they came from a quieter pass on the same
afternoon, so they are a fair timeline and are not comparable with anything.
Only a paired measurement is. Both tags, ten warm runs each, one sitting,
because chapter~\ref{ch:data-without-the-n-plus-1} learned the hard way that
measuring them hours apart measures the machine:

\begin{minted}[fontsize=\footnotesize]{text}
ch04  reviews as a list       5.743  5.924  5.946  6.130  6.480
                              6.560  6.662  6.730  7.074  7.260
ch05  reviews(first: 12)      8.333  8.385  8.858  9.017  9.260
                              9.368  9.556  9.691  9.838  10.178
\end{minted}

\index{pagination!honest cost}%
A mean of 9.2 against 6.5, for the same hundred
and twenty reviews. The connection is slower. A window function partitioned over
a parent key is more work than a plain scan, and at a hundred and twenty rows
the bound it buys is worth nothing at all, so this is what paying for insurance
looks like on a day nothing goes wrong.

I could have left this measurement out. Every argument in this section survives
without it, and it is the only number in the chapter that argues against the
change. It is here because a book that only prints the numbers that flatter its
decisions is not teaching anybody how to decide.

\subsection*{What the bound is actually for}

\index{cost analysis!listSize}%
Here is what was bought, and it is not really about pagination.

The connection publishes a directive the list never could:

\begin{graphqlsdl}
reviews(first: Int, after: String, last: Int, before: String): ReviewConnection!
  @listSize(
    assumedSize: 50
    slicingArguments: ["first", "last"]
    slicingArgumentDefaultValue: 10
    sizedFields: ["edges", "nodes"]
    requireOneSlicingArgument: false
  )
  @cost(weight: "10")
\end{graphqlsdl}

\code{slicingArguments} tells the cost analyzer that \code{first} and
\code{last} determine how big the answer is. Ask the analyzer what the catalog
query costs, varying only the page size, using the header
chapter~\ref{ch:hotchocolate-quickly} introduced:

\begin{minted}[fontsize=\footnotesize]{text}
reviews, as a list at tag ch04      fieldCost  30   typeCost    4
reviews(first: 2)                   fieldCost  41   typeCost    7
reviews(first: 5)                   fieldCost  71   typeCost   13
reviews(first: 50)                  fieldCost 521   typeCost  103
reviews, no page argument           fieldCost 121   typeCost   23
\end{minted}

Ten per review, exactly, with the no-argument case resolving to the declared
default of ten. And the list form, at the top: thirty, whatever came back.
Thirty for two reviews and thirty for two hundred thousand.

\index{cost analysis!unbounded lists}%
That is the argument, and it has almost nothing to do with clients wanting pages.
An unbounded list cannot be costed. If it cannot be costed it cannot be
budgeted, rate-limited or refused, so every protection
chapter~\ref{ch:security-hardening} builds on top of cost analysis is
protecting a number that was never true. The connection did not primarily make
the field pageable. It made the field's cost knowable, and under three
milliseconds on this data is a reasonable price for that.

\subsection*{The count that costs a round trip here and did not there}

\index{pagination!totalCount}%
One more measured surprise, and it contradicts something the last chapter said
plainly. Select \code{totalCount} on the nested connection and a fourth
statement appears:

\begin{minted}[fontsize=\footnotesize]{sql}
SELECT r.product_id AS "Key", count(*)::int AS "Count"
FROM reviews AS r
WHERE r.product_id = ANY (@productIds)
GROUP BY r.product_id
\end{minted}

Chapter~\ref{ch:data-without-the-n-plus-1} measured the opposite for
\code{browseProducts}, where the count arrived as a correlated subquery inside
the page's own statement and cost no extra round trip at all. Both measurements
are right, and the difference is structural rather than a caprice of the query
planner. A root connection is one page, so its count is one scalar and folds
into the same statement. A batched child connection is already a window over
groups, and the count of each group is a second aggregate over the same rows.

The practical form of that: on a root field, ask for \code{totalCount} freely.
On a batched child field, know that you are asking for a statement. Neither is a
rule about GraphQL, and both are things you find out by turning the database log
on, which is the only way anybody ever finds them out.

Three fields into the chapter, the schema can be read, paged and refetched.
Nothing in it has ever had to say no.
\index{pagination!nested connection|)}
